\documentclass[11pt]{letter}
\usepackage[T1]{fontenc}
\usepackage[utf8]{inputenc}
\usepackage{lmodern}
\usepackage[a4paper,margin=25mm]{geometry}
\usepackage{url}
\signature{Ku Kang\\CBRN Defense Research Institute, Korea Ministry of National Defense\\
Seoul, Republic of Korea\\\texttt{bisu9082@gmail.com}}
\address{}
\begin{document}
\begin{letter}{The Editors\\Atmospheric Environment}
\opening{Dear Editors,}

We submit the research paper \emph{``A global reanalysis reproduces when, but not where,
PM$_{10}$ is high across the dense monitoring network of an arid emirate''} by Jin Yoo, Jeongyun
Kim, Doo-Hee Lee and Ku Kang for consideration in \emph{Atmospheric Environment}.

Global reanalyses are increasingly used to assess exposure to coarse particulate matter where
monitors are few. Using the 19 regulatory monitors of the Emirate of Abu Dhabi, sixteen of which
share a cell of the CAMS reanalysis with another monitor, over 19 months, we show that the
reanalysis follows the day-to-day course of regional PM$_{10}$ but not the differences between
places: over the conurbation it combines a regional deficit on dusty days with a local offset of
about 60~$\mu$g\,m$^{-3}$ that does not change with regional dust, and up to half of the spatial
variance lies inside single model cells. We also show that the PM$_{10}$ values distributed for
this network by an open archive are 24-h running means, and that comparing them with
instantaneous model values manufactures an apparent benefit of statistical correction that
disappears when the two are averaged alike. We believe both results are of direct use to readers
who evaluate or apply reanalysis particulate matter in arid regions.

\textbf{Related manuscript.} A separate manuscript by the same authors is under consideration
in \emph{Atmospheric Environment} (ATMENV-D-26-01869, ``Meteorological covariates can mislead
zero-shot PM$_{2.5}$ nowcasting during Gulf dust extremes: an observation-validated benchmark'').
The two manuscripts share no data, model output or result. That study forecasts PM$_{2.5}$ from
the MERRA-2 aerosol reanalysis and ERA5 meteorology and validates it against OpenAQ monitors in
Manama, Kuwait City, Dhahran, Jeddah and Doha. The present study evaluates PM$_{10}$ from the
CAMS EAC4 reanalysis against the Abu Dhabi regulatory network, none of whose monitors appears in
the other manuscript, and involves no forecasting.

This manuscript has not been published and is not under consideration elsewhere. All authors
have approved the submission. All data, code and result files are openly available at
\url{https://github.com/bisu9082/serra_camscast}. The authors declare no competing interests.

\closing{Sincerely,}
\end{letter}
\end{document}
